\chapter{Comparing and viewing structures}
\label{ch:compare}

Three tools help you look beyond a single file. The Homology Searcher finds
proteins related to yours. The Structure Aligner lays two or more structures on
top of each other and measures how different they are. The Structure Viewer
shows any of your structures in 3D in a web browser. None of them is needed to
prepare a simulation, but they help you choose and check your starting
structure.

\begin{glance}
\glanceitem{What it does}{BLAST search for related proteins (3a);
  superposition with RMSD and TM-score (3b); interactive 3D view (ua).}
\glanceitem{You need it when}{You look for a better template or a homolog's
  AlphaFold model, compare a crystal structure with an NMR or predicted one, or
  want to see what a step did.}
\glanceitem{Before you start}{A loaded structure. Internet access for BLAST.
  A web browser on the same computer for the viewer.}
\glanceitem{Where to find it}{\menupath{3a}{Homology Searcher},
  \menupath{3b}{Structure Aligner}, \menupath{ua}{Structure Viewer}; on the
  command line \texttt{-{}-pdbview}.}
\glanceitem{You get}{\file{blast_export.xml}, AlphaFold models in
  \file{homologs/}, aligned files such as \file{1AKK_aligned.pdb}, saved viewer
  scenes.}
\end{glance}

\section{Finding related proteins with BLAST}
\label{sec:compare-blast}

\term{BLAST} compares the amino-acid sequence of your protein with all the
sequences in a database and lists the most similar ones. ProPrep sends the
search to the public BLAST server of the NCBI (the US National Center for
Biotechnology Information), so you do not need BLAST on your computer. You need
an internet connection, and patience.

For each hit, BLAST reports the \term{identity} (the percentage of aligned
positions with the same amino acid) and the \term{E-value}: the number of hits
this good you would expect by pure chance. The smaller the E-value, the more
certain the relationship; values below about $10^{-5}$ are usually taken as
real relatives.

\subsection{Step by step: a BLAST search for hemoglobin alpha}

\begin{steps}
  \item Type \key{3a}, then \key{1} (\prompt{Run BLAST homology search}).
  \item If several structures are loaded, ProPrep shows them in a table. Type
        the number of the one whose sequence you want:
\begin{screen}
                   Available Structures in Workspace
┃ Option ┃ Type                ┃ File Path                         ┃      Size ┃
│   1    │ Biological Assembly │ .../asm/1HHO_biological_assembly… │ 396.31 KB │
│   2    │ RCSB PDB            │ .../asm/4HHB.pdb                  │ 462.74 KB │
Select structure to use (1): §ans«2»
\end{screen}
  \item Choose where the sequence comes from. \key{3} takes it from the loaded
        structure; \key{1} lets you paste a sequence; \key{2} reads a FASTA
        file. Then choose the chain:
\begin{screen}
Sequence input methods:
1. Enter sequence directly
2. Load from FASTA file
3. Use sequence from loaded structure (4HHB)
Choose input method (1): §ans«3»

Available chains in structure 4HHB:
Chain A: 141 residues
Chain B: 146 residues
Chain C: 141 residues
Chain D: 146 residues
Select chain (A): §ans«A»
Using sequence from chain A (141 residues)
\end{screen}
  \item Choose the program and the database. For a protein, the program is
        \key{2} (\prompt{blastp}, protein against protein). For the database,
        type \key{3} (swissprot):
\begin{screen}
Database Options:
1. nr - Non-redundant protein sequences
2. refseq_protein - NCBI protein reference sequences
3. swissprot - UniProtKB/Swiss-Prot
4. pdb - Protein Data Bank proteins

Note: Options 1 and 3 allow retrieval of AlphaFold structures for homologs.
Select database (1): §ans«3»
\end{screen}
        \prompt{swissprot} holds about half a million carefully checked
        proteins; \prompt{pdb} holds only proteins with an experimental
        structure. Both are far smaller and faster than \prompt{nr}, the
        default, which holds hundreds of millions of sequences. Choose
        \prompt{swissprot} if you want AlphaFold models of the hits, and
        \prompt{pdb} if you want experimental structures.
  \item Press \key{Enter} to keep the E-value threshold of 10.0. Type a
        small maximum number of hits, for example \typed{20}.
  \item Choose the \term{scoring matrix}, the table that says how similar two
        amino acids are. The default, \prompt{1. BLOSUM62 - General purpose},
        suits almost every search; the screen explains the others.
  \item Check the summary and type \key{y}:
\begin{screen}
Ready to run BLAST with the following parameters:
  Program: blastp
  Database: swissprot
  Sequence length: 141 characters
  E-value threshold: 10.0
  Maximum hits: 20
Run BLAST with these parameters? [y/n] (y): §ans«y»
This may take several minutes depending on sequence length and server load.
\end{screen}
\end{steps}

\begin{watchout}
\textbf{A BLAST search can take hours.} The search waits in the NCBI's public
queue, and ProPrep shows no progress and offers no way to cancel. In our test,
after 10 minutes a warning appeared:
\begin{screen}
BiopythonWarning: BLAST request G3XE7CNX016 is taking longer than 10 minutes,
consider re-issuing it
\end{screen}
You can ignore it and keep waiting. Our search finished after almost four hours
(\prompt{BLAST search completed in 13741s}). The only way to stop earlier is
\key{Ctrl}+\key{C}. To keep waits short, use \prompt{swissprot} or
\prompt{pdb}, not \prompt{nr}, and ask for few hits.
\end{watchout}

\subsection{Reading the results}

When the search is done, ProPrep saves the full result as
\file{blast_results_blastp.xml} and shows a table (shortened here):
\begin{screen}
BLAST search completed in 13741s. Found 20 hits.

===== BLAST Results Summary =====
Query: unnamed protein product
Database: swissprot
Number of hits: 20

┃ Hit ┃ Description                                  ┃ Score ┃ E-value  ┃ Identity ┃
│ 1   │ sp|P69905.2| RecName: Full=Hemoglobin sub... │ 284   │ 2.28e-99 │ 100.0%   │
│ 2   │ sp|P01923.1| RecName: Full=Hemoglobin sub... │ 283   │ 9.61e-99 │ 99.3%    │
│ 3   │ sp|Q9TS35.2| RecName: Full=Hemoglobin sub... │ 279   │ 1.84e-97 │ 98.6%    │
...
│ 20  │ sp|Q7M3B6.1| RecName: Full=Hemoglobin sub... │ 271   │ 3.04e-94 │ 95.7%    │

View detailed alignment for specific hits? [y/n] (y): §ans«y»
Enter hit numbers (comma-separated) (1): §ans«1»
\end{screen}
All 20 hits are alpha chains of hemoglobin from different animals. The
best hit is the human alpha chain itself (UniProt \pdb{P69905}), with 100\% identity, as
expected for a search with a human protein. The detailed view shows the
sequence alignment:
\begin{screen}
Score: 726.0 bits (284.263)
E-value: 2.28437e-99
Identities: 141/141 (100.0%)
Positives: 141/141 (100.0%)
Query range: 1-141
Subject range: 2-142

Alignment:
Query     1  VLSPADKTNVKAAWGKVGAHAGEYGAEALERMFLSFPTTKTYFPHFDLSHGSAQVKGHGK  60
             VLSPADKTNVKAAWGKVGAHAGEYGAEALERMFLSFPTTKTYFPHFDLSHGSAQVKGHGK
Sbjct     2  VLSPADKTNVKAAWGKVGAHAGEYGAEALERMFLSFPTTKTYFPHFDLSHGSAQVKGHGK  61
...
\end{screen}
The query starts at 1 and the database entry (\prompt{Sbjct}) at 2, because
the UniProt sequence includes the first methionine, which is removed in the
mature protein and absent from the crystal structure.

Option \key{2} of the Homology Searcher menu (\prompt{View BLAST results})
shows the results again later.

\subsection{AlphaFold models of the hits}

Option \key{3} (\prompt{Retrieve AlphaFold structures for hits}) downloads the
predicted structures of chosen hits. It works only after a search of
\prompt{nr} or \prompt{swissprot}. ProPrep lists the hits with their UniProt
IDs and asks how to choose:
\begin{screen}
Selection Methods:
1. Top N hits
2. By E-value threshold
3. By identity % threshold
4. Manual selection (enter ranks)
5. All hits
Select method (1): §ans«1»
Enter number of top hits to retrieve (max 20) (5): §ans«2»
Saving structures to: .../asm/homologs
...
✓ Successfully retrieved 2 AlphaFold structures
\end{screen}
The models are saved in the folder \file{homologs/}, for example as
\file{AF-P01923-F1-model_v6.cif}. ProPrep then offers to make one of
them the structure for later steps
(\prompt{Would you like to select a homolog for downstream processing?}); we
chose the second, and ProPrep said
\prompt{Selected P01923 for downstream processing}.

\subsection{Saving the results}

Option \key{4} (\prompt{Export alignments and annotations}) saves the results
as \file{blast_export.xml}, in the standard BLAST XML format that other
programs read. If you answer \key{y} to the next question, it also saves a
table that opens in a spreadsheet:
\file{blast_export_summary.tsv}.

\subsection{Homology models}

\begin{needsoftware}
Option \key{5} (\prompt{Build homology model (MODELLER)}) builds a model of
your protein from the structure of a relative, a so-called
\term{homology model}. It needs MODELLER with a valid license key
(Section~\ref{sec:cleaning-fixer}). Without it the menu shows
\prompt{MODELLER not installed} and the option is blocked. It was not tried for
this manual.
\end{needsoftware}

\section{Superimposing structures with the Structure Aligner}
\label{sec:compare-align}

To compare two structures of the same or related proteins, you first have to
lay one on top of the other as closely as possible. This is a
\term{superposition}. One structure stays where it is: the
\term{reference}. The others, the \term{targets}, are rotated and moved onto it.
Then ProPrep measures what is left over.

Two numbers describe the result:
\begin{description}
  \item[RMSD] The \term{root-mean-square deviation}: the typical distance
        between matching atoms after superposition, in Å. Near 0 means
        practically identical. In our tests, the same protein measured by X-ray
        and by NMR differed by about 1.7~Å, and a crystal structure and its
        AlphaFold prediction by 0.34~Å. A few large local differences (a loop,
        a chain end) can raise the RMSD a lot.
  \item[TM-score] A score from 0 to 1 that says how similar the overall
        \emph{fold} is, and is less affected by small floppy regions. Above
        about 0.5 two proteins have the same fold; close to 1 they are
        practically identical.
\end{description}

You also choose which atoms to compare: the \term{backbone} (the atoms N, CA,
C and O that every amino acid has; the default), only the \term{alpha carbons}
(CA, one point per residue), or all atoms. Backbone or CA comparisons work even
when the side chains differ.

\subsection{Step by step: crystal against NMR structure of cytochrome c}

The example compares horse cytochrome~c measured by X-ray (\pdb{1HRC}) with
the same protein measured by NMR (\pdb{1AKK}). Both were loaded from the RCSB
first; \pdb{1AKK} was loaded last.

\begin{steps}
  \item Type \key{3b}, then \key{1} (\prompt{Align multiple PDB structures}).
        Step~1 adds the latest RCSB structure by itself and lists other PDB
        files in the project folder:
\begin{screen}
Step 1: Load Structures
✓ Selected RCSB PDB: .../align/1AKK.pdb
✓ Added 1 workspace structure(s)
...
Available local structure files:
  • 1HRC.pdb

Options:
  1. Browse for local structure file (PDB/CIF)
  2. Download PDB from database
  3. Done (proceed with alignment)
Select option (1): §ans«1»
\end{screen}
  \item Add the second structure: \key{1} opens the file browser (choose
        \file{1HRC.pdb}); \key{2} downloads one by PDB code. Then type \key{3}
        (Done). For each structure, press \key{Enter} at
        \prompt{Use entire structure from \ldots ? [y/n] (y)}; answering
        \key{n} lets you choose a model and chains.
  \item Step~2: choose the reference. ProPrep advises an experimental
        structure if you want to keep its ligands, or an AlphaFold model if you
        want a complete chain. Type \key{2} (\file{1HRC.pdb}).
  \item Step~3: press \key{y} at \prompt{Use all common residues for
        alignment?}.
  \item Step~4: choose how matching residues are found:
\begin{screen}
Step 4: Residue Mapping Strategy
  [1] Sequence-based (automatic) - for homologous proteins with similar
      sequences
  [2] Structure-based (TM-align) - for proteins with same fold but different
      sequences
  [3] Manual specification - for identical numbering or custom residue selection
Select method (1): §ans«1»
\end{screen}
  \item Step~5: press \key{Enter} for backbone atoms (mode \key{2}).
  \item ProPrep aligns the sequences (\prompt{Alignment: 104/104 identical
        (100.0\%)}) and asks whether to match non-protein groups such as the
        heme by their position
        (\prompt{Interactively map non-protein residues using spatial
        proximity? [y/n] (n)}). Press \key{Enter} (no) unless you want to
        carry ligands from one structure to another.
  \item Step~6: press \key{Enter} at \prompt{Save intermediate alignment
        steps for debugging?}.
\end{steps}

ProPrep adds residues one by one and shows how the RMSD develops, then the
result:
\begin{screen}
RMSD Progress: 1AKK
...
  Final RMSD: 1.679 Å

Saving Aligned Structures...
✓ 1AKK_aligned.pdb
✓ 1HRC_aligned.pdb
...
Select Target Structure for Downstream Processing
...
│ 1 │ 1AKK_aligned.pdb │ 1AKK.pdb │
Select target structure (or 'n' to skip) (1): §ans«1»
...
✓ Structure alignment complete!
\end{screen}
It takes a few seconds. Both structures are saved in the common position:
\file{1AKK_aligned.pdb} and \file{1HRC_aligned.pdb}. An RMSD of 1.7~Å between
a crystal and an NMR structure of the same protein is normal: NMR structures
are more spread out, especially at the surface.

The last question makes one aligned target the structure for later steps.
Answer \key{n} if you only wanted the comparison.

\begin{watchout}
An aligned file that you select here outranks a fresh download in later tools
(Section~\ref{sec:structure-switching}). If you align only to compare, answer
\key{n} at \prompt{Select target structure}, or do the comparison in a separate
project.
\end{watchout}

\subsection{Structure-based alignment with TM-align}

When two proteins have the same shape but different sequences, sequence
matching fails. Strategy \key{2} uses \term{TM-align}, which matches residues
by their position in space. The example compared hemoglobin alpha in the
crystal (\pdb{4HHB}, chain A) with its AlphaFold prediction (\pdb{P69905}).

The AlphaFold download is an mmCIF file, which the aligner's browser does not
list (Section~\ref{sec:structure-alphafold}). So we downloaded the model
in PDB format from the AlphaFold Database website, put it into the project
folder as \file{AF-P69905-F1-model_v6.pdb}, and picked that. For \pdb{4HHB} we answered \key{n} at
\prompt{Use entire structure} and chose chain A only. With the AlphaFold model
as reference and strategy \key{2}:
\begin{screen}
Performing TM-align structure-based alignment...

Aligning: 4HHB.pdb
  Aligning all chains combined
  Reference: 142 CA atoms
  Target: 141 CA atoms

  TM-score (normalized by target): 0.9943
  TM-score (normalized by reference): 0.9874
  RMSD: 0.335 Å
\end{screen}
A TM-score of 0.99 and an RMSD of 0.34~Å mean the prediction and the crystal
structure are practically identical. The score is given twice because it can
be scaled by the length of either structure. TM-align always uses the alpha
carbons, so there is no atom-mode question. If a structure has several chains,
ProPrep first asks whether to align all chains together or one chain.

\subsection{Other choices in the aligner}

\begin{itemize}
  \item \textbf{Part of the protein only.} Answer \key{n} at
        \prompt{Use all common residues for alignment?} and choose \key{1} to
        type ranges such as \typed{A:1-40}. With residues 1--40, strategy
        \key{3} (manual, same numbering) and CA atoms, \pdb{1HRC} and
        \pdb{1AKK} gave 1.991~Å. A wrongly typed range is reported
        (\prompt{Invalid format. Use chain:start-end}) and skipped without
        asking again.
  \item \textbf{Repeated runs.} A second run does not overwrite: it writes
        \file{1HRC_aligned_2.pdb} and so on. At the start ProPrep asks
        \prompt{Include these aligned structures for further alignment?};
        answer \key{y} to build on them.
  \item \textbf{Aligning on redox sites} (option \key{2}) needs sites found by
        the Redox Site Detector \key{2a} first (Chapter~\ref{ch:redox}).
  \item \textbf{Viewing the result} (option \key{3}) prints the RMSD after each
        step of a sequence-based run.
\end{itemize}

\begin{watchout}
In version 1.23.0:
\begin{itemize}
  \item Option \key{3} (\prompt{View alignment results}) after a TM-align run
        prints \prompt{Error executing Structure Aligner.\_view\_alignment\_results:
        'step'} and a long error message. ProPrep keeps running; the aligned
        files are fine.
  \item The browser says \prompt{PDB/CIF} but refuses \file{.cif} files.
  \item If you answer \key{n} to \prompt{Include these aligned structures for
        further alignment?}, the loaded structures are not offered either
        (\prompt{Current structures: 0}); add them with the browser.
  \item With fewer than two structures ProPrep stops with
        \prompt{At least 2 structures required for alignment}.
  \item After any alignment, option \key{2} looks available. Without detected
        sites it says \prompt{Please run redox site detection first (PDB Filter
        module)}; the right tool is \key{2a}.
\end{itemize}
\end{watchout}

\section{Viewing structures in 3D}
\label{sec:compare-viewer}

ProPrep has a 3D viewer that runs in your web browser. ProPrep starts a small
web server on your own computer and gives you its address, for example
\url{http://localhost:8765/viewer}. ``\term{localhost}'' means ``this
computer'': the page is not on the internet and nobody else can see it.

\subsection{Step by step: opening the viewer}

\begin{steps}
  \item Type \key{ua} (Structure Viewer), then \key{1}
        (\prompt{Launch interactive viewer}).
  \item Choose one or more structures from the table, by number, range
        (\typed{1-3}), list (\typed{1,2}) or \typed{all}:
\begin{screen}
Select one or more structures to view in the interactive viewer.
Structures will be displayed together (superimposed if aligned).
...
┃ Option ┃ Type              ┃ File Path                          ┃
│   1    │ Aligned Target    │ .../align/1AKK_aligned.pdb         │
│   2    │ Aligned Reference │ .../align/1HRC_aligned_2.pdb       │
│   3    │ RCSB PDB          │ .../align/4HHB.pdb                 │
│   4    │ AlphaFold         │ .../align/AF-P69905-F1-model_v6.…  │
...
Select structure(s) (1): §ans«1,2»
✓ Selected 2 structures
✓ Structure viewer launched at http://localhost:8765/viewer (2 structures)
\end{screen}
  \item ProPrep tries to open the page in your default browser. If no page
        appears, open your browser yourself and type the address from the last
        line.
\end{steps}

The page is called ``ProPrep Interactive Structure Viewer''. It shows the
protein as a cartoon and ligands as balls and sticks. When you use
\texttt{proprep-web} (Chapter~\ref{ch:starting}), the viewer appears in a panel
next to the ProPrep terminal instead.

Notes on the address:
\begin{itemize}
  \item The number after \texttt{localhost:} is the \term{port}. ProPrep uses
        8765 if it is free, otherwise the next free number (8766, 8767,
        \ldots). Always use the address ProPrep prints.
  \item The viewer runs as long as ProPrep runs. There is no ``stop'' option;
        it stops when you leave ProPrep with \key{x}.
  \item If ProPrep runs on a remote computer that you reach by SSH, it cannot
        open a browser there. It then prints a suggested command to make the
        address reachable from your own computer (not tried for this manual).
\end{itemize}

\begin{tip}
Many tools open the viewer by themselves: viewing metadata
(Section~\ref{sec:structure-metadata}), \prompt{View local files together},
the PDB Filter's viewer question, the protonation rationale
(Section~\ref{sec:protonation-propka}) and the orientator. Read the printed
address each time: a second viewer can open on a new port. In our test the
electron-density option opened port 8766 while 8765 was still running.
\end{tip}

\begin{watchout}
Choose structures from the same alignment run. The table lists the latest
``Aligned Target'' and ``Aligned Reference'', which can come from
\emph{different} runs: here \file{1AKK_aligned.pdb} came from the first run and
\file{1HRC_aligned_2.pdb} from a later one, so they are not in the same frame.
Select the target at the end of each alignment, or use a fresh project.
\end{watchout}

\subsection{Other viewer options}

\begin{description}[style=nextline]
  \item[2 View the loaded MD trajectory] Shows a simulation as a movie. It
        needs a topology and trajectory loaded with \key{1a} option \key{2}
        (Chapter~\ref{ch:analysis}); it was not tried for this chapter.
  \item[3 Show available annotations] Lists what ProPrep can mark on the
        structure: redox sites, titratable residues, filter selections, bond
        classifications and aligned structures. In our test it said
        \prompt{No annotations found in workspace} even though aligned
        structures existed.
  \item[4 Save the scene] Saves what the open viewer shows (structures and
        display styles) in a file ending in \file{.scene.json}. The suggested
        name contains the date and time.
  \item[5 Load a saved scene] Lists the saved scenes and applies one to the
        open viewer.
  \item[6 Electron density] Shows the experimental electron density of an
        X-ray structure around the model (next section).
\end{description}

\subsection{Electron density}

An X-ray structure is a model built into a map of \term{electron density}.
Looking at the map shows how well each part of the model is supported by the
data. Option \key{6} fetches two maps from the European PDB (PDBe): the
\term{2mFo-DFc} map, the density the model was built into, and the
\term{mFo-DFc} difference map, which shows where data and model disagree.
ProPrep's own explanation adds: \prompt{Trust the difference map more}.

You choose a structure and a margin around the model (default 5~Å). For the
high-resolution crambin structure \pdb{3NIR}:
\begin{screen}
3NIR.pdb: fetching the maps of 3NIR from PDBe...
✓ 3NIR.pdb: density of 3NIR (2mFo-DFc sigma 1.033, mFo-DFc difference sigma
0.263). Its atoms sit at +10.17 sigma of the 2mFo-DFc map on average
✓ Structure viewer launched at http://localhost:8765/viewer (1 structure)
\end{screen}
A model in its own density usually sits around $+3$~sigma; much more means a
very clear map. The maps are saved in the project folder (\file{3nir_2fofc.ccp4}
and others). They can be large: about 5.5~MB each as downloaded, and 22.9~MB
each after cutting around the model.

Maps exist only for X-ray entries deposited with their measured data, and only
for structures still in their original position (as downloaded, filtered or
protonated, but not superimposed onto another structure). For \pdb{4HHB}
ProPrep said \prompt{PDBe has no maps for 4HHB} and opened the viewer without
density. NMR and cryo-EM entries have no such maps.

\subsection{A quick look from the command line: \texttt{-{}-pdbview}}

To look at a structure without starting a project, start ProPrep like this:
\begin{terminal}
proprep --full-menu --pdbview 1UBQ
\end{terminal}
ProPrep downloads \pdb{1UBQ}, opens the viewer and prints the address:
\begin{screen}
✓ Structure viewer launched at http://localhost:8766/viewer (1 structure)
\end{screen}
You can also give a file: \texttt{proprep -{}-full-menu -{}-pdbview 1UBQ.pdb}.
ProPrep then shows the main menu with the structure loaded. Nothing is
recorded in a session file. Without \texttt{-{}-full-menu} (or
\texttt{-{}-workflow}) ProPrep first asks which menu mode you want. A file that
does not exist gives \prompt{File not found} and the main menu.

\section{If something goes wrong}
\label{sec:compare-problems}

\begin{description}[style=nextline]
  \item[BLAST runs for a very long time] This is the NCBI queue. Wait, or stop
        with \key{Ctrl}+\key{C} and search again with \prompt{swissprot} or
        \prompt{pdb} and fewer hits.
  \item[``Retrieve AlphaFold structures'' finds nothing] It works only after a
        search of \prompt{nr} or \prompt{swissprot}.
  \item[The aligner does not list my AlphaFold file] It lists only
        \file{.pdb} files. Download the PDB-format model from the AlphaFold
        Database website.
  \item[``Error executing Structure Aligner.\_view\_alignment\_results'']
        Viewing results after a TM-align run fails in 1.23.0. The aligned
        files are fine; open them in the viewer.
  \item[The viewer page does not open] Open the printed address in your
        browser yourself. Use exactly the port number ProPrep printed.
  \item[The viewer address no longer works] ProPrep has been closed; the
        viewer stops with it. Start it again from \key{ua}.
  \item[Two aligned structures do not overlap in the viewer] They come from
        different alignment runs. Align again and select the target.
\end{description}

See also Chapters~\ref{ch:troubleshooting} and~\ref{ch:knownissues}.
